\chapter*{บรรณานุกรม}
\addcontentsline{toc}{chapter}{บรรณานุกรม}

\begin{enumerate}[label={[\arabic*]}, leftmargin=2em, itemsep=6pt]
    \item ชีวะ ทัศนา, และคณะ. (2567). \textit{การประยุกต์ใช้อากาศยานไร้คนขับและปัญญาประดิษฐ์ฝังตัวสำหรับการประเมินดัชนีพืชพรรณในสวนผลไม้เศรษฐกิจภาคตะวันออก}. วารสารวิทยาศาสตร์และเทคโนโลยี มหาวิทยาลัยราชภัฏรำไพพรรณี, \textit{18}(2), หน้า 45--58.
    \item ชีวะ ทัศนา. (2568). \textit{ฟิสิกส์เชิงคำนวณและปัญญาประดิษฐ์ทางการเกษตรดิจิทัล}. เอกสารคำสอนวิชาฟิสิกส์ประยุกต์ขั้นสูง, คณะวิทยาศาสตร์และเทคโนโลยี มหาวิทยาลัยราชภัฏรำไพพรรณี.
    \item Gitelson, A. A., Kaufman, Y. J., Stark, R., \& Rundquist, D. (2002). Novel algorithms for remote estimation of vegetation fraction. \textit{Remote Sensing of Environment}, \textit{80}(1), 76--87.
    \item Louhaichi, M., Borman, M. M., \& Johnson, D. E. (2001). Spatially located platform and aerial photography for documentation of grazing impacts on wheat. \textit{Geocarto International}, \textit{16}(1), 65--70.
    \item Woebbecke, D. M., Meyer, G. E., Von Bargen, K., \& Mortensen, D. A. (1995). Color indices for weed identification under various soil, residue, and lighting conditions. \textit{Transactions of the ASAE}, \textit{38}(1), 259--269.
    \item Hunt, E. R., Cavigelli, M., Daughtry, C. S., McMurtry, J., \& Walthall, C. L. (2005). Evaluation of digital photography from model aircraft for remote sensing of crop biomass and nitrogen status. \textit{Precision Agriculture}, \textit{6}(4), 359--378.
    \item Thassana, C. (2025). Spectral analysis of vegetation indices using edge mobile computing for orchard management. \textit{Journal of Applied Agricultural Physics}, \textit{12}(3), 112--126.
    \item Google Inc. (2026). \textit{Flutter documentation: Building high performance cross-platform applications}. Retrieved from \url{https://docs.flutter.dev}
\end{enumerate}

\clearpage
